\ficha{Operação, 1907 a 1909}

\bloco{Fontes lidas}

Não há volume parlamentar para esses anos. Li por inteiro quatro peças de
imprensa: a carta de Serzedelo Corrêa ao \textit{Diario de Noticias}, de
Lisboa, reproduzida em \textit{O Paiz} de 27 de março de 1907; a nota de
\textit{O Paiz} de 24 de junho de 1908 sobre o parecer do orçamento da receita
para 1909, de que Serzedelo era relator; a entrevista de Serzedelo ao
\textit{Correio da Manhã} de 27 de maio de 1909; e duas peças de \textit{O
Paiz} com voz do jornal, de 21 de março de 1907 e de 1º de dezembro de 1909.

\bloco{Síntese}

Com a Caixa em funcionamento, o debate se desloca do padrão para a operação:
o que ela estabilizou, o que fazer com o papel inconversível que continua em
circulação, a que taxa cobrar os direitos em ouro na alfândega e a quem
pertence o ouro depositado. Serzedelo, vencido em 1906, passa a reconhecer o
resultado sem abandonar a exigência de resgate.

\bloco{Serzedelo Corrêa depois da lei}

Em Lisboa, na carta ao \textit{Diario de Noticias}, ele explica a Caixa ao
público português:

\cit[I:per178691_1907_08210:3]{O beneficio que a caixa de conversão está
destinada a produzir é o da estabilização relativa do
cambio}{\jor{opaiz19070327}{O Paiz}{27 mar. 1907}{3}}

e acrescenta o que em 1906 negava ao projeto original:
\tr[I:per178691_1907_08210:3]{Para evitar a baixa, a caixa tem os recursos
accumulados na alta}. A inversão é coerente com a declaração de 29 de novembro
de 1906 (ficha 1), mas a carta atribui à Caixa as operações de câmbio que o
Senado havia separado dela. É apresentação para o estrangeiro, não análise.

No parecer da receita, segundo o resumo de \textit{O Paiz}, o relator
\tr[I:per178691_1908_08665:1]{affirma que a caixa produziu a estabilização das
taxas cambiaes} e aconselha o Congresso a \tr[I:per178691_1908_08665:1]{apparelhar
o governo com os recursos para o resgate do papel-moeda} (\jor{opaiz19080624}{O
Paiz}{24 jun. 1908}{1}). A avaliação positiva vem com a mesma ressalva de 1906:
a emissão da Caixa deve substituir papel, não se somar a ele.

A entrevista de 1909 trata da revisão da tarifa aduaneira, isto é, da taxa a
que se cobram os direitos em ouro, e não do padrão monetário. É nesse contexto
que ele diz: \tr[I:per089842_1909_02872:1]{A tendencia geral é para a
metallização da nossa moeda} (\jor{cm19090527}{Correio da Manhã}{27 maio
1909}{1}), propondo que a totalidade dos direitos de importação seja paga em
ouro para compensar a perda de receita.

\bloco{O depósito da Caixa: Campista e a resposta de O Paiz}

Em entrevista ao \textit{Le Brésil}, em Paris, já fora do governo, David
Campista sugere tomar o depósito da Caixa como encaixe metálico de parte do
papel-moeda. \textit{O Paiz} transcreve o telegrama do \textit{Jornal do
Commercio} com a proposta, que começa por \tr[I:per178691_1909_09189:1]{Os dez
milhões do deposito actual}, e responde em nota de redação, em primeira pessoa
do plural. Para o jornal, Campista esquece

\cit[I:per178691_1909_09189:1]{que o deposito actual de dez milhões de libras
responde pela emissão feita}{\jor{opaiz19091201}{O Paiz}{1 dez. 1909}{1}}

É o argumento do depósito como garantia dos portadores de bilhetes, que
reaparece em 1910 e em 1914 (fichas 5 e 7).

\bloco{A carteira de câmbio e a taxa de 15}

Em março de 1907, uma coluna satírica de \textit{O Paiz}, sem assinatura
visível no OCR, trata da saída de Custódio Coelho, diretor da carteira de
câmbios do Banco do Brasil. A tese da coluna é que o governo pediu ao Banco que
segurasse o câmbio em 15 e que este

\cit[I:per178691_1907_08204:1]{começou a recalcar o cambio, cuja tendencia para
a alta o proprio Sr. Campista}{\jor{opaiz19070321}{O Paiz}{21 mar. 1907}{1}}

reconhecera. A mesma coluna diz que Custódio, adversário da quebra do padrão,
acabou aceitando a prefixação da taxa de 15. A frase é do colunista, em tom
irônico, e não pode ser usada como declaração de Custódio. Não a transcrevo
porque o OCR do trecho ficou abaixo do limiar da conferência mecânica.

\bloco{Achados}

\begin{enumerate}[leftmargin=*,nosep]
\item Serzedelo muda de avaliação sobre o aparelho entre 1906 e 1907, e a
mudança acompanha as emendas que ele mesmo pedira. O que não muda é a
exigência de resgate do papel inconversível.
\item A entrevista de 1909 pertence ao debate da tarifa em ouro, não ao do
padrão. Citar a frase sobre a \textit{metallização} fora desse contexto
distorce o sentido.
\item O argumento de que o ouro da Caixa responde pelos bilhetes emitidos
aparece em 1909 na voz de \textit{O Paiz}, antes de aparecer na de qualquer
parlamentar lido.
\end{enumerate}

\bloco{Limites}

Tudo nesta ficha vem de imprensa lida no OCR, sem conferência na imagem, salvo
as datas das edições. O parecer da receita é conhecido só pelo resumo do
jornal. O telegrama com a entrevista de Campista é transcrição de terceira mão.
